\begin{claim}{C2.2}{Proved}{Comparison with McHugh and Angella}
The specialization of McHugh's Theorem 6.3, Proposition 6.4, equation
(7), and Section 7.1 to \eqref{eq:C-hodge-input} agrees with C2.1.
The corresponding relations and Figure 5 in Angella's Section 4 also
agree.  No correction to these specific formulas is required.
\end{claim}
\begin{proof}
McHugh's quotient $K^{p,q}$ is zero in intermediate total degrees
on every arrow and square just used: a closed corner of an arrow is
the image of the incoming differential and a closed corner of a square
is in both images.  Thus $k^{p,q}=0$ for $0<p+q<6$.
Also $E_2^{3,1}=0$ by C1.1 and $H^{1,1}(X;\R)=0$ because
$H^2(X;\R)=0$.  Reading the formulas in
\citet[Theorem 6.3, Proposition 6.4, (7), pp. 20--25]{McHugh2018}
therefore gives
\[
h_{BC}^{1,1}=2h^{0,1}=2,\qquad
h_{BC}^{1,2}=h^{1,2}=h,\qquad
h_{BC}^{2,2}=-1+(h+1)+h=2h.
\]
The remaining entries in his table are precisely the boundary entries
of \eqref{eq:C-BC-A}.  In
\citet[Section 4, pp. 13--16]{Angella2018}, the relation
\[
h_{BC}^{2,2}=2h_{BC}^{2,1}-2h^{0,1}+2
\]
gives $2h$, and the boundary values reduce to those above because
$h^{2,0}=h^{0,2}=0$.  This comparison concerns exactly the formulas
used here; it does not certify every assertion in either article.
\end{proof}
